\scanpage{78}
ne change pas quand on intervertit les roles de $g_{i_0}$ et $f$, car
\[
V/g_{i_0}^*(\mathcal J_f)\simeq A/(\mathcal J_{g_{i_0}}+\mathcal J_f)
\]
-- car $\nu$ est la longueur du schéma artinien local, intersection des \inserted{ss-schémas} $f(S)$ et $g_{i_0}(S)$ de $X$ -- on peut aussi l'interpréter comme une multiplicité d'intersection. Ceci posé, Je présume que l'hom
\[
T_\ell\longrightarrow T_\ell^I/T_\ell
\]
est le produit par $\mu$ de l'injection canonique $T_\ell-T_\ell^I$, correspondant à l'indice $i_0$. Il faudrait que je demande à Carlos de me le confirmer.

Reste le cas $\nu=1$, qui semble demander un traitement séparé.

\unsure{Le|\deleted{\illegible{1}}Le} \deleted{\illegible{1}}fait à vérifier (pour les groupes fondamentaux premi\unsure{e\uline{rs}|\uline{ers}}\margin{\deleted{\unsure{?!}}\\?!} à $p$) c'est que \unsure{le|l'} hom. ext. $\pi_1(\eta)\longrightarrow\pi_1(U)\simeq\pi_1(U_s)$ est égal à $\deleted{\unsure{\mu}}\unsure{k|\kappa}_{i_0}\circ(\mu.\operatorname{id}_T)$, où $\unsure{k|\kappa}_{i_0}$ est
\margin{Ceci doit être indépendant de la valeur de $\nu$ !}
\margin{77}
